\section{Reinforcements, Replacements, Transfers, and Upgrading}\label{sec:13.0}

\ruleslabel{General Rule}

\begin{generalrule}
Both Players receive reinforcements and replacements during the course of the game. In addition, both Players may transfer units at pre-determined times. Similarly, the Axis Player is permitted to upgrade certain units at specific times during the course of the game.
\end{generalrule}

\ruleslabel{Procedure}

\begin{procedure}
A. Axis reinforcements and replacements appear on the map during the Axis Reinforcement Phase of the indicated Game-Turn. In addition, unit transfer and upgrading takes place at this time. The Axis Player may place reinforcements and replacements anywhere within a Zone that the unit would normally be allowed to move into according to the restrictions of Case 6.6 (even in a Display whose corresponding box, triangle, or circle is currently occupied by an Enemy unit).
\end{procedure}

\textbf{B.} Yugoslav reinforcements and replacements appear on the map at the
beginning of the Yugoslav Movement Phase, except for variable guerrilla
reinforcements, which appear in the Guerrilla Reinforcement Phase (see 7.4). In
addition, transfer of Soviet and pro-Yugoslav Bulgarian units takes place at the
beginning of the Yugoslav Movement Phase. Yugoslav units are placed on the map
according to the requirements of Cases 13.2, 13.4, and 13.6.

\ruleslabel{Cases}

\subsection{Axis Reinforcements}\label{sec:13.1}

The Axis Player is due reinforcements as described in the Reinforcement, Transfer, and Upgrading Schedule (13.9).

\subsection{Yugoslav Reinforcements}\label{sec:13.2}

\case{13.21} The Yugoslav Player is due predetermined Partisan reinforcements (including Tito) on Game-Turn 2 (see 13.92).

\case{13.22} The Yugoslav Player is due Soviet and Bulgarian reinforcements on Game-Turn 15 (see 13.92). In addition, there are three optional Soviet reinforcing units that are available to the Yugoslav Player at this time. However, if these units are used, the Yugoslav Player loses 15 Victory Points in the ensuing Victory Point Stage (see 14.3).

\subsection{Axis Reinforcements Triggered By Events}\label{sec:13.3}

The Axis Player is due reinforcements as certain events are triggered during the course of the game. These are:

\textbf{A.} Italian Withdrawal (see 10.2); B. Italian Surrender (see 10.3); C. Tito is located (see 9.42);

\textbf{D.} The Yugoslav Player first accumulates 45 Victory Points or the first Yugoslav brigade is placed on the map.

Note: The units available as a result of these events are listed in Case 13.91.

\subsection{Yugoslav Variable Guerrilla Reinforcements}\label{sec:13.4}

The Yugoslav Player is due variable guerrilla reinforcements during all Guerrilla Reinforcement Phases (see 7.4).

\subsection{Axis Replacements}\label{sec:13.5}

Axis replacements must be drawn from units that have been previously eliminated from play. Replacements may never be accumulated; if they are not used or not available on a given Game-Turn, they are permanently lost.

\case{13.51} Starting with the Axis Reinforcement Phase of Game-Turn 2, the Axis Player receives one Serbian unit (Axis Player’s choice) as a replacement each Game-Turn.

\case{13.52} Starting with the Axis Reinforcement Phase of Game-Turn 3, the Axis Player receives two Croat units (Axis Player’s choice) as replacements each Game-Turn.

\subsection{Yugoslav Guerrilla Replacements}\label{sec:13.6}

\case{13.61} If, at the beginning of any Yugoslav Movement Phase of Game-Turn 3 (or after), there are no Partisan units on the map, the Yugoslav Player is immediately eligible for Partisan replacements (see 13.63).

\case{13.62} If, at the beginning of any Yugoslav Movement Phase between Game-Turns 2 and 8 (inclusive), there are no Chetnik units of any allegiance on the map, the Yugoslav Player is immediately eligible for pro-Yugoslav Chetnik replacements (see 13.63).

\case{13.63} If the Yugoslav Player is due replacements, he rolls a single die for Partisans and/or a single die for Chetniks. The result(s) is the number of replacement groups of the appropriate type immediately available to the Yugoslav Player.

\case{13.64} Partisan replacements may be placed in the Hide-away circle of the
Bosnia, Croatia, Serbia, Dalmatia, and/or Slovenia Zone Display.

Chetnik replacements (always placed on the map as pro-Yugoslav) may be placed in the Hide-away circle of the Montenegro and/or Serbia Zone Display.

\subsection{Transferring Units}\label{sec:13.7}

\case{13.71} The Reinforcement, Transfer, and Upgrading Schedule (13.9) will ask the Axis Player to transfer (remove) specific units on the map at various times during the game. In order to transfer a unit, the Axis Player simply picks it up and removes it from the map during the Axis Reinforcement Phase. A unit that is transferred may not return to play unless specified by 13.9.

\case{13.72} Transfer of Axis units is not mandatory. However, during the Victory Point Stage of each Game-Turn in which there are Axis units on the map whose transfer had previously been called for, the Yugoslav Player receives 5 Victory Points for each such unit (see 14.2).

\case{13.73} If an Axis unit has been eliminated when its transfer is called for, this transfer is ignored.

\case{13.74} The Yugoslav Player may transfer only Soviet and pro-Yugoslav Bulgarian units after Game-Turn 15. This is performed exactly like Axis transfer, except it takes place at the beginning of the Yugoslav Movement Phase. The Yugoslav Player loses 10 Victory Points for each Soviet and pro-Yugoslav Bulgarian unit on the map during the Victory Point Stage of Game-Turns 16 and 17 (see 14.3).

\subsection{Upgrading Axis Units}\label{sec:13.8}

\case{13.81} The Reinforcement, Transfer, and Upgrading Schedule (13.9) will list the times at which specific Axis units will be upgraded. In order to upgrade a unit, the Axis Player simply flips it over, revealing a new unit with a stronger Combat Strength.

\case{13.82} If a unit that is to be upgraded is eliminated from play, its upgrading is ignored. Note: Eliminated Croat units that re-deploy onto the map as Axis replacements (see 13.52) after their scheduled upgrading are automatically upgraded.

\subsection{Reinforcement, Transfer, And Upgrading Schedule}\label{sec:13.9}

All reinforcements, transfers, and units to be upgraded are given by Game-Turn. Each unit’s designation is listed, followed by a letter and number in parentheses which indicate nationality and Combat Strength. Abbreviations are as follows: G = German; C = Croat; P = Partisan; B = Bulgarian; S = Soviet. All units listed are reinforcements unless otherwise noted.

\case{13.91} \textbf{Axis}

\textbf{At Start}

\textbf{A.} Place in Box \#6 of the Allied Progress Track (available after
Italian Withdrawal): 100L(G10), 173(G6), 297(G12), 373(G10), 92(G6).

\textbf{B.} Place in Box \#7 of the Allied Progress Track (available after
Italian Surrender): 181(G18), 264(G15), 371(G15), 1Cos(G9).

\textbf{C.} Place aside (available after locating Tito): 501SS(G1).

\textbf{D.} Place aside (available after the Yugoslav Player first accumulates
45 Victory Points or the first Yugoslav brigade is placed on the map): 25(B12),
27(B12). These reinforcements must be deployed (and remain) in Serbia, only in
an Objective Display labeled ``Bulgarian Occupation.''

\textbf{Game-Turn 2:} 342(G18), 125(G3), Gd(C3), 1Ust(C2), 2Ust(C2), 3Ust(C2).

\textbf{Game-Turn 3:} 113(G18), 4Ust(C2), 5Ust(C2). \emph{Transfer:} 125(G3).

\textbf{Game-Turn 4:} 1(C1), 2(C1), 3(C1), 4(C1), 5(C1), 1Mt(C2), 2Mt(C2),
3Mt(C2), 4Mt(C2), 5Mt(C2). \emph{Transfer:} 113(G18).

\textbf{Game-Turn 5:} \emph{Transfer:} 342(G18).

\textbf{Game-Turn 6:} 7SS(G16).

\textbf{Game-Turn 8:} 369(G10), 6(C1). \emph{Upgrade:} 704(G12), 714(G12),
717(G12), 718(G12). \emph{Transfer:} 117L(G15).

\textbf{Game-Turn 9:} 1(G20).

\textbf{Game-Turn 10:} 1Pz(G30). \emph{Transfer:} 104L(G15), 1Pz(G30).
Employment of 1Pz during Game-Turn 10 and each succeeding Game-Turn
automatically yields 5 Victory Points to the Yugoslav Player.

\textbf{Game-Turn 12:} 13SS(G8), 187(G9).

\textbf{Game-Turn 13:} 202Pz(G10).

\textbf{Game-Turn 14:} 21SS(G8), 23SS(G4), 98(G20), 392(G10).

\textbf{Game-Turn 15:} 4SS(G20), 22(G24), 11LW(G12), 104L(G15), 117L(G15).
These units must be placed anywhere within Macedonia. \emph{Upgrade:} all Croat
units. \emph{Transfer:} all Bulgarian units (may not be delayed by the Axis
Player).

\case{13.92} \textbf{Yugoslav}

\textbf{Game-Turn 2:} 10 groups (each P1), Tito unidentified: Serbia
Yugoslav/Partisan Hide-away circle; 4 groups (each P1): Slovenia
Yugoslav/Partisan Hide-away circle; 3 groups (each P1): Montenegro
Yugoslav/Partisan Hide-away circle.

\textbf{Game-Turn 15:} 75(S30), 68(S30), 64(S30), 10Gd(S33), 31Gd(S33),
4Gd(S42), 5(S12), 1(B15), 2(B15), 3(B15), 4(B15), 5(B15), 7(B15), 8(B15),
6(B18), 9(B18). Place in any Yugoslav/Partisan box, triangle, or circle in
Serbia or Macedonia (Soviet units), or simply Macedonia (Bulgarian units).
None of these units may leave their original Zone of placement on Game-Turn 15.

The following optional Soviet reinforcements are also available at this time:
18Gd(S33), 20Gd(S33), 7(S38). For each one employed, the Yugoslav Player loses
5 Victory Points in the ensuing Victory Point Stage.
